% !PRINTED 112
\weberpage{112}

\begin{center}
\textbf{\S\ 34. \gesperrt{Die Funktionen $\eta(\omega)$, $f(\omega)$, $f_1(\omega)$, $f_2(\omega)$.}}
\end{center}

Es sind bereits im § 24 die Funktionen $\eta(\omega)$, $f(\omega)$, $f_1(\omega)$,
$f_2(\omega)$ erwähnt, die sich dort als einwertige Funktionen von $\omega$ bei der
Darstellung der $\vartheta$-Funktionen durch unendliche Produkte fast von selbst
einstellten, die sich aber nicht ergaben bei der zweiten Darstellung durch unendliche
Reihen. Damit im Zusammenhange steht ein bemerkenswerter Umstand, daß viele unserer
Formeln, z.\,B. die zur Transformation zweiter Ordnung gehörigen (17), § 32 oder die
Formel (23), § 32, leicht verifiziert werden können durch die unendlichen Produkte,
dagegen schwer oder gar nicht durch die unendlichen Reihen. Darum war es für uns von
Interesse, diese Resultate ohne die Benutzung des einen oder anderen dieser Ausdrücke
aus der Transformationstheorie herzuleiten, und ebenso sollen nun auch aus dieser
Quelle die Funktionen $\eta(\omega)$, $f(\omega)$, $f_1(\omega)$, $f_2(\omega)$ und
ihre Grundeigenschaften gewonnen werden.

Die Formel (6) des vorigen Paragraphen ergibt für $n=3$:
\[
\sqrt{3}\,\vartheta_{11}(3u,\,3\omega)\,\vartheta_{00}\,\vartheta_{10}\,\vartheta_{01}
= 2\,\vartheta_{11}\!\left(\tfrac{1}{3}\right)\vartheta_{11}(u)\,
\vartheta_{11}\!\left(\tfrac{1}{3}+u\right)\vartheta_{11}\!\left(\tfrac{1}{3}-u\right),
\]
und wenn man differentiiert und dann $u=0$ setzt nach (4):
\[
3\sqrt{3}\,\vartheta'_{11}(0,\,3\omega) = 2\pi\,\vartheta_{11}\!\left(\tfrac{1}{3}\right)^3,
\]
oder indem man $\omega$ durch $\dfrac{\omega}{3}$ ersetzt,
\[
3\sqrt{3}\,\vartheta'_{11} = 2\pi\left[\vartheta_{11}\!\left(\tfrac{1}{3},\,\tfrac{\omega}{3}\right)\right]^3.
\]
Setzt man also
\weq{1}{\eta(\omega) = \frac{1}{\sqrt{3}}\,\vartheta_{11}\!\left(\frac{1}{3},\,\frac{\omega}{3}\right),}
so folgt in der Bezeichnung übereinstimmend mit § 24 (9):
\weq{2}{\vartheta'_{11} = \pi\,\vartheta_{00}\,\vartheta_{01}\,\vartheta_{10} = 2\pi\,\eta(\omega)^3}
und aus § 25, (4) erhält man für $\eta(\omega)$ die Reihenentwickelung
\weq{3}{\eta(\omega) = q^{\frac{1}{12}} \sum_{-\infty,\,\infty}^{\nu} (-1)^{\nu}\, q^{3\nu^2+\nu}.}

Die Funktion $\eta(\omega)$ ist für ein rein imaginäres $\omega$ (reelles $q$) \gesperrt{reell},
und für ein unendlich großes $\omega$ (d.\,h. verschwindendes $q$) ist
\[
q^{-\frac{1}{12}}\,\eta(\omega) = 1.
\]
